\section{Critical Discussion of the Results}

Overall, the results support three main conclusions. First, porting the \emph{SafeSwitch} logic was successful: in the artificial scenarios the system reproduces the qualitative behavior expected from the reference work and clearly differs from a \emph{naive} fallback in the most critical cases. Second, migrating the cellular component to Simu5G does not alter the causal chain underlying the architecture: interference on 802.11p, cellular handovers, and gap variations continue to be reflected in platoon behavior in a way that is consistent with what is discussed in \cite{segata2023multi-technology}. Third, exact numerical agreement with the paper should not be expected, because the transition from SimuLTE to Simu5G inevitably changes the radio stack, the available physical models, and their related default values.

The most important point of the comparison is therefore not to establish whether 5G is better than LTE in absolute terms within this experiment, but to verify whether the studied mechanism remains meaningful and transferable in an updated environment. From this perspective, the outcome is positive: no evident functional regressions emerge, the system continues to react plausibly to connectivity degradation, and the realistic scenario preserves the key phenomena highlighted by the paper, especially the link between packet losses, gap increase, and controller state changes.

At the same time, the differences observed in FER peaks and maximum inter-vehicle distances require caution in quantitative interpretation. To make the comparison tighter, a systematic calibration campaign on radio parameters, handovers, and propagation models would be needed, together with identical replication conditions between the LTE and 5G stacks. The contribution of this thesis should therefore be read primarily as a validation of the architecture and experimental setup, rather than as a definitive measurement of the advantage of one cellular technology over the other.
